
本附录介绍\BookRecordIndex{T0319}索伯列夫空间，以及\BookRecordIndex{T0334}伪全纯曲线理论所需的\BookRecordIndex{T0220}拉普拉斯算子 $L^p$ \BookRecordIndex{T0408}正则性理论。
在最后一节，我们仅假设二维\BookRecordIndex{T0220}拉普拉斯算子的内部\BookRecordIndex{T0408}正则性，便证明具有全实边界条件的\BookRecordIndex{T0334}伪全纯曲线的基本\BookRecordIndex{T0408}正则性与\BookRecordIndex{T0192}紧性定理。

第一节阐述\BookRecordIndex{T0319}索伯列夫空间，包括基本定义、\BookRecordIndex{T0035}Lipschitz 域上光滑函数在\BookRecordIndex{T0319}索伯列夫空间中的稠密性证明、\BookRecordIndex{T0261}庞加莱不等式及\BookRecordIndex{T0318}索伯列夫嵌入定理的证明，以及一个\BookRecordIndex{T0174}迹定理的版本。B.2 节依照 Gilbarg 与 Trudinger~\srcite{140} 的书证明 \BookRecordIndex{T0012}Calderón--Zygmund 不等式；B.3 节讨论\BookRecordIndex{T0220}拉普拉斯算子的内部\BookRecordIndex{T0408}正则性；B.4 节解释\BookRecordIndex{T0334}伪全纯曲线的\BookRecordIndex{T0330}椭圆自举技术。前三节虽然在 B.4 节的应用中只需处理 $n=2$，但仍对 $n$ 个变量的函数展开论述，因为在这种更一般的情形下，基本理论并不比二维情形更难。

\section{\BookRecordIndex{T0319}索伯列夫空间}\label{sec:B-1}

设 $\Omega\subset\mathbb R^n$ 是开集。记 $C^\infty(\overline\Omega)$ 为 $\mathbb R^n$ 上光滑函数限制到 $\overline\Omega$ 所得的空间，$C_0^\infty(\Omega)$ 为 $\Omega$ 上紧支撑的光滑函数空间。先提及对光滑紧支撑函数 $u,v:\mathbb R^n\to\mathbb R$ 成立的两个基本不等式：当 $1/p+1/q=1$ 时的 Hölder 不等式
\[
  \|uv\|_{L^1}\le\|u\|_{L^p}\|v\|_{L^q},
\]
以及 Young 不等式
\[
  \|u*v\|_{L^p}\le\|u\|_{L^1}\|v\|_{L^p}.
\]
这里
\[
  u*v(x)=\int_{\mathbb R^n}u(x-y)v(y)\,dy
\]
表示\BookRecordIndex{T0205}卷积，而
\[
  \|u\|_{L^p}=\left(\int_{\mathbb R^n}|u|^p\right)^{1/p}
\]
表示 $1\le p<\infty$ 时的 $L^p$ 范数。

\paragraph{\BookRecordIndex{T0293}弱导数。} 设 $u:\Omega\to\mathbb R$ 局部可积，固定多重指标 $\nu=(\nu_1,\ldots,\nu_n)$。若局部可积函数 $u_\nu:\Omega\to\mathbb R$ 满足
\[
  \int_\Omega u(x)\partial^\nu\phi(x)\,dx
  =(-1)^{|\nu|}\int_\Omega u_\nu(x)\phi(x)\,dx
\]
对所有 $\phi\in C_0^\infty(\Omega)$ 成立，则称 $u_\nu$ 是 $u$ 对应于 $\nu$ 的\emph{\BookRecordIndex{T0293}弱导数}。\BookRecordIndex{T0293}弱导数若存在，便由 $u$ 几乎处处唯一确定，记为
\[
  \partial^\nu u:=u_\nu.
\]
数 $|\nu|:=\nu_1+\cdots+\nu_n$ 称为导数的阶。由散度定理，任意 $C^k$ 函数 $u:\Omega\to\mathbb R$ 都有不超过 $k$ 阶的\BookRecordIndex{T0293}弱导数，且它们与通常的导数一致。

现在固定非负整数 $k$ 和 $1\le p\le\infty$。\emph{\BookRecordIndex{T0319}索伯列夫空间} $W^{k,p}(\Omega)$ 定义为所有满足下述条件的 $L^p(\Omega)$ 函数 $u$（的等价类）所成的空间：对每个 $|\nu|\le k$，\BookRecordIndex{T0293}弱导数 $\partial^\nu u$ 存在且属于 $L^p$。当 $1\le p<\infty$ 时，$u\in W^{k,p}(\Omega)$ 的 $W^{k,p}$ 范数定义为
\[
  \|u\|_{k,p}=\left(\int_\Omega\sum_{|\nu|\le k}|\partial^\nu u(x)|^p\,dx\right)^{1/p}.
\]
当 $p=\infty$ 时，$W^{k,\infty}$ 范数定义为所有 $|\nu|\le k$ 的\BookRecordIndex{T0293}弱导数 $\partial^\nu u$ 的 $L^\infty$ 范数之最大值。空间 $W^{k,p}_{\mathrm{loc}}(\Omega)$ 由局部 $p$ 可积函数 $u:\Omega\to\mathbb R$（的等价类）组成，它们限制到 $\Omega$ 中任何预紧开子集 $Q$ 后都属于 $W^{k,p}(Q)$。空间 $W^{k,p}_0(\Omega)$ 定义为 $C_0^\infty(\Omega)$ 在 $W^{k,p}(\Omega)$ 中的闭包，因此它是 $C_0^\infty(\Omega)$ 关于 $W^{k,p}$ 范数的完备化。

\begin{exercise}\label{ex:B-1-1}
证明 $W^{k,p}(\Omega)$ 是巴拿赫空间，并且在 $1<p<\infty$ 时自反。证明 $W^{k,p}_0(\Omega)$ 在 $1\le p<\infty$ 时\BookRecordIndex{T0213}可分。提示：对适当整数 $N$，把 $W^{k,p}(\Omega)$ 看作 $L^p(\Omega,\mathbb R^N)$ 的闭子空间。巴拿赫空间的任意闭子空间完备；自反巴拿赫空间的任意闭子空间自反。证明\BookRecordIndex{T0213}可分性时，可利用如下事实：每个光滑函数 $u:\Omega\to\mathbb R$ 都可由有理系数多项式序列逼近，而且各阶导数在每个紧集上一致收敛。
\end{exercise}
\begin{exercise}\label{ex:B-1-2}
设 $\Omega\subset\mathbb R^n$ 是有界开集。
\begin{enumerate}[label=(\roman*)]
\item 证明 $u\in W^{k+1,p}(\Omega)$ 当且仅当 $u\in W^{1,p}(\Omega)$，且\BookRecordIndex{T0293}弱导数 $\partial_i u=\partial u/\partial x_i$ 对 $i=1,\ldots,n$ 都属于 $W^{k,p}(\Omega)$。
\item 若 $u\in W^{k,p}(\Omega)$ 且 $v\in W^{k,\infty}(\Omega)$，证明 $uv\in W^{k,p}(\Omega)$，并且
\[
  \|uv\|_{k,p}\le c\|u\|_{k,p}\|v\|_{k,\infty},
\]
其中常数 $c$ 只依赖于 $k,n$。提示：证明比看上去困难。难点是证明：若 $u\in W^{1,p}(\Omega)$ 且 $v\in W^{1,\infty}(\Omega)$，则 $uv$ 的\BookRecordIndex{T0293}弱导数遵循 Leibniz 法则
$\partial_i(uv)=(\partial_i u)v+u(\partial_i v)$，$i=1,\ldots,n$。这些函数显然属于 $L^p$；有了这一步，结论就可归纳得到。证明 Leibniz 法则需要下面的命题~\ref{prop:B-1-4}。先在 $u$ 光滑时证明，再在 $\Omega$ 的紧子集上以一列光滑函数逼近 $u$。
\end{enumerate}
\end{exercise}

\paragraph{光滑函数逼近。} 接下来的目标是证明：对一大类域 $\Omega\subset\mathbb R^n$，\BookRecordIndex{T0319}索伯列夫空间 $W^{k,p}(\Omega)$ 可视为 $C^\infty(\overline\Omega)$ 关于 $W^{k,p}$ 范数的完备化。若开集 $\Omega\subset\mathbb R^n$ 的边界局部可表示为某个 Lipschitz 函数的图像，就称它为\emph{\BookRecordIndex{T0035}Lipschitz 域}。明确地说，对每个 $x\in\partial\Omega$，须有 $x$ 的邻域 $U$、非零向量 $\xi\in\mathbb R^n$、常数 $\delta>0$，以及满足 $f(0)=0$、$|\eta|<\delta$ 时 $|f(\eta)|<\delta$ 的 Lipschitz 连续函数 $f:\xi^\perp\to\mathbb R$，使得
\[
  \Omega\cap U=\{x+\eta+t\xi\mid\eta\in\xi^\perp,\ |\eta|<\delta,\ f(\eta)<t<\delta\}.
\]
这个定义的动机在于下述观察。

\begin{exercise}\label{ex:B-1-3}
设 $1\le p<\infty$，$\Omega\subset\mathbb R^n$ 为有界 \BookRecordIndex{T0035}Lipschitz 域。证明：对每个 $u\in W^{k,p}(\Omega)$，存在一列开集 $\Omega_j\subset\mathbb R^n$ 以及函数 $u_j\in W^{k,p}(\Omega_j)$，使得
\[
  \overline\Omega\subset\Omega_j,\qquad
  \lim_{j\to\infty}\|u_j-u\|_{W^{k,p}(\Omega)}=0.
\]
提示：先设 $u\in W^{k,p}(\Omega)$ 支撑于某边界点的足够小的邻域 $U$。证明函数 $u_\varepsilon(x):=u(x+\varepsilon\xi)$ 在 $\Omega$ 上按 $W^{k,p}$ 范数收敛到 $u$。一般情形则以适当的单位分解乘 $u$，并用习题~\ref{ex:B-1-2}(ii)。
\end{exercise}

\begin{proposition}\label{prop:B-1-4}
若 $\Omega\subset\mathbb R^n$ 是有界 \BookRecordIndex{T0035}Lipschitz 域，则对 $1\le p<\infty$ 和任意非负整数 $k$，$C^\infty(\overline\Omega)$ 在 $W^{k,p}(\Omega)$ 中稠密。
\end{proposition}
\begin{proof}
取 $u\in W^{k,p}(\Omega)$。先设 $u$ 在 $\Omega$ 中有紧支撑。取光滑非负函数 $\rho:\mathbb R^n\to\mathbb R$，使
\[
  \operatorname{supp}\rho\subset B_1,\qquad\int_{\mathbb R^n}\rho(x)\,dx=1,
\]
并对 $\delta>0$ 定义 $\rho_\delta(x)=\delta^{-n}\rho(\delta^{-1}x)$。于是 $\rho_\delta*u$ 光滑且有紧支撑。由\BookRecordIndex{T0293}弱导数定义直接得到，$\rho_\delta*u$ 的通常导数为
\begin{equation}\label{eq:B-1-1}
  \partial^\nu(\rho_\delta*u)=\rho_\delta*\partial^\nu u.
\end{equation}
故由 Young 不等式，
\begin{equation}\label{eq:B-1-2}
  \|\rho_\delta*u\|_{W^{k,p}}\le\|u\|_{W^{k,p}}.
\end{equation}
容易验证，对每个紧支撑连续函数 $f:\Omega\to\mathbb R$，$\delta\to0$ 时 $\rho_\delta*f$ 一致收敛到 $f$。由于紧支撑连续函数在 $L^p(\Omega)$ 中稠密（$1\le p<\infty$），取 $k=0$ 的不等式~\eqref{eq:B-1-2} 表明对任意 $f\in L^p(\Omega)$，$\rho_\delta*f$ 都按 $L^p$ 范数收敛到 $f$。于是由~\eqref{eq:B-1-1}，$\rho_\delta*u$ 按 $W^{k,p}$ 范数收敛到 $u$。这证明每个紧支撑的 $W^{k,p}$ 函数都可由紧支撑光滑函数序列逼近。

现在考虑一般情形，令 $\varepsilon>0$。由习题~\ref{ex:B-1-3}，有开集 $\Omega'\subset\mathbb R^n$ 和函数 $u'\in W^{k,p}(\Omega')$，使 $\overline\Omega\subset\Omega'$ 且 $\|u'-u\|_{W^{k,p}(\Omega)}\le\varepsilon/2$。不失一般性，可设 $u'$ 在 $\Omega'$ 中有紧支撑。由刚证明的结论，存在 $\Omega'$ 上的紧支撑光滑函数 $v$，使 $\|v-u'\|_{W^{k,p}(\Omega')}\le\varepsilon/2$。命题得证。
\end{proof}

命题~\ref{prop:B-1-4} 的证明中那样的一列光滑函数 $\rho_\delta$ 称为\emph{$\delta$ 函数的逼近}，算子 $u\mapsto\rho_\delta*u$ 称为\emph{\BookRecordIndex{T0151}光滑化算子}。

\begin{exercise}\label{ex:B-1-5}
本习题说明\BookRecordIndex{T0319}索伯列夫空间 $W^{k,p}$ 在右侧与 $C^{k-1,1}$ 微分同胚复合后不变；这里 $C^{k-1,1}$ 微分同胚是指其 $(k-1)$ 阶导数 Lipschitz 连续的 $C^{k-1}$ 微分同胚。设 $\Omega,\Omega'\subset\mathbb R^n$ 是有界 \BookRecordIndex{T0035}Lipschitz 域，$\psi:\Omega'\to\Omega$ 是一个 $C^{k-1,1}$ 微分同胚（即它是两个适当的闭包邻域之间的 $C^{k-1,1}$ 微分同胚的限制）。证明：若 $u\in W^{k,p}(\Omega)$，则 $u\circ\psi\in W^{k,p}(\Omega')$，且
\[
  \|u\circ\psi\|_{W^{k,p}(\Omega')}\le c\|u\|_{W^{k,p}(\Omega)},
\]
其中 $c$ 与 $u$ 无关。提示：利用习题~\ref{ex:B-1-2} 和命题~\ref{prop:B-1-4}，对 $k$ 作归纳。
\end{exercise}

\paragraph{\BookRecordIndex{T0261}庞加莱不等式。} 若 $u\in W^{1,p}(\Omega)$ 的所有导数都为零，它必在 $\Omega$ 的每个连通分支上为常数，这一点并非显然；证明需要下面的基本估计。一如既往，$\nabla u$ 表示 $u$ 的梯度。回忆函数的平均值，是其在定义域上的积分除以该域的体积。

\begin{lemma}[\BookRecordIndex{T0261}庞加莱不等式]\label{lem:B-1-6}
设 $1\le p<\infty$，$\Omega\subset\mathbb R^n$ 是有界开域。则对 $u\in W_0^{1,p}(\Omega)$，有
\[
  \|u\|_{L^p(\Omega)}\le\operatorname{diam}(\Omega)\|\nabla u\|_{L^p(\Omega)}.
\]
若 $\Omega=Q^n=(0,1)^n$ 是单位方块，则每个满足 $\int_{Q^n}u=0$ 的 $u\in W^{1,p}(\Omega)$ 都满足
\[
  \|u\|_{L^p(Q^n)}\le n\|\nabla u\|_{L^p(Q^n)}.
\]
\end{lemma}
\begin{proof}
第一项只须对 $u\in C_0^\infty(\Omega)$ 证明。不失一般性，设 $\Omega\subset\{x_n>0\}$ 且 $0\in\partial\Omega$。于是
\[
  u(x)=\int_0^{x_n}\partial_n u(x_1,\ldots,x_{n-1},t)\,dt.
\]
由于 $|x_n|\le\operatorname{diam}(\Omega)$，Hölder 不等式给出
\[
  |u(x)|^p\le\operatorname{diam}(\Omega)^{p-1}
  \int_0^\infty|\partial_nu(x_1,\ldots,x_{n-1},t)|^p\,dt.
\]
将两边在 $\mathbb R^{n-1}\times[0,\operatorname{diam}(\Omega)]$ 上积分，得到
\[
  \int_\Omega|u|^p\le\operatorname{diam}(\Omega)^p\int_\Omega|\partial_nu|^p.
\]
第一项得证。第二项对 $n$ 作归纳。$n=1$ 时估计是一个简单练习。设已对 $n\ge1$ 成立，并取平均值为零的 $u\in C^\infty(\overline{Q^{n+1}})$。定义
\[
  v(t)=\int_{Q^n}u(x_1,\ldots,x_n,t)\,dx_1\cdots dx_n.
\]
由于 $v\in C^\infty(\overline{Q^1})$ 且平均值为零，
\[
  \int_0^1|v(t)|^p\,dt\le\int_0^1|\dot v(t)|^p\,dt
  \le\int_{Q^{n+1}}|\partial_{n+1}u|^p\,dx.
\]
最后一步来自 Hölder 不等式。
对每个 $t$，归纳假设给出
\[
  \int_{Q^n}|u(x,t)-v(t)|^p\,dx
  \le n^p\int_{Q^n}|\nabla u(x,t)|^p\,dx.
\]
再对 $t$ 积分，得到 $\|u-v\|_{L^p}\le n\|\nabla u\|_{L^p}$。结合前一估计便有 $\|u\|_{L^p(Q^{n+1})}\le(n+1)\|\nabla u\|_{L^p(Q^{n+1})}$。这证明了光滑函数的第二项；一般情形由命题~\ref{prop:B-1-4} 得出。
\end{proof}

\begin{corollary}\label{cor:B-1-7}
设 $\Omega\subset\mathbb R^n$ 是有界开域，$u\in W^{1,p}(\Omega)$，且 $\partial u/\partial x_j\equiv0$，$j=1,\ldots,n$。则 $u$ 在 $\Omega$ 的每个连通分支上为常数。若还满足 $u\in W_0^{1,p}(\Omega)$，则 $u\equiv0$。
\end{corollary}
\begin{proof}
先设 $\Omega$ 是方块且 $u$ 的平均值为零。取平均值为零的光滑函数序列 $u_\nu\in C^\infty(\overline\Omega)$，使之按 $W^{1,p}$ 范数逼近 $u$。\BookRecordIndex{T0261}庞加莱不等式表明 $u_\nu$ 按 $L^p$ 范数收敛到零，故 $u=0$。于是，若 $u$ 的一阶导数在某方块上为零，它在那里便为常数；从而 $\nabla u\equiv0$ 的 $u\in W^{1,p}(\Omega)$ 在 $\Omega$ 的每个连通分支上为常数。若 $u\in W_0^{1,p}(\Omega)$，取 $u_\nu\in C_0^\infty(\Omega)$ 按 $W^{1,p}$ 范数收敛到 $u$。则 $\|\nabla u_\nu\|_{L^p}\to0$，故由引理~\ref{lem:B-1-6}，$\|u_\nu\|_{L^p}\to0$，从而 $u\equiv0$。
\end{proof}

上述推论也可由下一习题\BookRecordIndex{T0326}推出：它说明对任意开集 $\Omega\subset\mathbb R^n$，\BookRecordIndex{T0319}索伯列夫空间 $W_{\mathrm{loc}}^{1,\infty}(\Omega)$ 可自然地与 $\Omega$ 上局部 Lipschitz 连续函数空间 $C_{\mathrm{loc}}^{0,1}(\Omega)$ 等同。

\begin{exercise}\label{ex:B-1-8}
\begin{enumerate}[label=(\roman*)]
\item 若 $u\in L^1_{\mathrm{loc}}(\Omega)$，证明 $u_\delta(x)=\rho_\delta*u(x)$ 对几乎每个 $x\in\Omega$ 都收敛到 $u(x)$（见命题~\ref{prop:B-1-4} 的证明）。
\item 证明 $C_{\mathrm{loc}}^{0,1}(\Omega)\subset W_{\mathrm{loc}}^{1,\infty}(\Omega)$。提示：设 $u:\Omega\to\mathbb R$ 局部 Lipschitz 连续，固定 $\xi\in\mathbb R^n$。证明序列 $u_j(x)=j\bigl(u(x)-u(x-\xi/j)\bigr)$\sourcefootnote{原书差商括号内印为 $u(x-\xi/j)-u(x)$，其极限是沿 $\xi$ 方向导数的相反数，与紧随其后的\BookRecordIndex{T0293}弱导数恒等式符号相反；按该恒等式调整差商次序。} 有子序列在每个紧子集 $K\subset\Omega$ 上的 $L^2(K)$ 中弱收敛。证明极限函数 $u^\xi:\Omega\to\mathbb R$ 是 $u$ 沿 $\xi$ 方向的\BookRecordIndex{T0293}弱导数，即对所有 $\phi\in C_0^\infty(\Omega)$，
\[
  \int_\Omega u^\xi\phi=-\int_\Omega u\langle\nabla\phi,\xi\rangle.
\]
证明 $u^\xi$ 局部有界。
\item 证明 $W_{\mathrm{loc}}^{1,\infty}(\Omega)\subset C_{\mathrm{loc}}^{0,1}(\Omega)$。提示：若 $u\in W_{\mathrm{loc}}^{1,\infty}(\Omega)$，证明对 $\delta>0$，函数 $u_\delta=\rho_\delta*u$ 在 $B_r(x)$ 上局部 Lipschitz 连续，其 Lipschitz 常数可取 $c=\sup_{B_{r+\delta}(x)}|\nabla u|$。再利用 (i) 证明 $u$ 局部 Lipschitz 连续（可能需要在零测集上重新定义）。
\end{enumerate}
\end{exercise}

\paragraph{延拓。} 下一个命题是一个有用的辅助结果：只要 $\Omega$ 有 $C^{k-1,1}$ 边界，每个 $W^{k,p}(\Omega)$ 函数都可延拓为较大域上紧支撑的 $W^{k,p}$ 函数，而且延拓后函数的范数可由原函数范数控制。

\begin{proposition}\label{prop:B-1-9}
设 $\Omega\subset\mathbb R^n$ 是边界为 $C^{k-1,1}$ 类的有界开集，$\Omega'\subset\mathbb R^n$ 是满足 $\overline\Omega\subset\Omega'$ 的开集。\sourcefootnote{原书在此未明言 $k\ge1$，但 $C^{k-1,1}$ 边界与证明中的 $k$ 项反射公式在 $k=0$ 时均无定义；命题应至少按 $k\ge1$ 理解。}则存在有界线性算子
\[
E:W^{k,p}(\Omega)\to W_0^{k,p}(\Omega'),
\]
使得 $Eu|_\Omega=u$ 对每个 $u\in W^{k,p}(\Omega)$ 成立。
\end{proposition}
\begin{proof}
证明取自~\srcite{140}。先考虑 $\Omega=\mathbb H^n=\{x_n>0\}$ 的情形。在 $x_n\le0$ 时以
\[
(E_0u)(x_1,\ldots,x_{n-1},x_n):=\sum_{i=1}^kc_i u(x_1,\ldots,x_{n-1},-x_n/i)
\]
定义 $E_0:C^{k-1,1}(\mathbb H^n)\to C^{k-1,1}(\mathbb R^n)$，其中 $c_1,\ldots,c_k$ 满足
\[
\sum_{i=1}^kc_i\left(-\frac1i\right)^m=1,\qquad m=0,\ldots,k-1.
\]
容易验证，不超过 $k-1$ 阶的导数在边界上匹配；若 $|x|\ge R$ 时 $u(x)=0$，则 $|x|\ge kR$ 时 $E_0u(x)=0$；对紧支撑函数还有估计
\[
\|E_0u\|_{W^{k,p}(\mathbb R^n)}\le c_0\|u\|_{W^{k,p}(\mathbb H^n)}.
\]
现在设 $\Omega$ 是边界为 $C^{k-1,1}$ 类的任意有界开集，并记 $B_r:=\{x\in\mathbb R^n\mid |x|<r\}$。可取开覆盖 $\overline\Omega\subset U_0\cup\cdots\cup U_N$、开集 $U_j'$，以及 $C^{k-1,1}$ 微分同胚 $\psi_j:U_j'\to B_k$，$j=1,\ldots,N$，使得 $\overline U_0\subset\Omega$、$\overline U_j\subset U_j'\subset\Omega'$，且
\[
\psi_j(U_j)=B_1,\qquad\psi_j(U_j'\cap\Omega)=B_k\cap\mathbb H^n
\]
对 $j=1,\ldots,N$ 成立。取单位分解 $\beta_j:\mathbb R^n\to[0,1]$，满足 $\operatorname{supp}\beta_j\subset U_j$，且求和上界为 $N$。\sourcefootnote{原书此处求和上界印为 $n$，但单位分解与下式均以 $j=0,\ldots,N$ 为指标；据此改为 $N$。}
\[
\sum_{j=0}^{N}\beta_j(x)=1\qquad(x\in\Omega).
\]
定义 $E:C^{k-1,1}(\Omega)\to C_0^{k-1,1}(\Omega')$ 为
\[
Eu=\beta_0u+\sum_{j=1}^N\bigl(E_0((\beta_ju)\circ\psi_j^{-1})\bigr)\circ\psi_j.
\]
由命题~\ref{prop:B-1-4} 和习题~\ref{ex:B-1-5}，$E$ 延拓为从 $W^{k,p}(\Omega)$ 到 $W_0^{k,p}(\Omega')$ 的有界线性算子。
\end{proof}

\paragraph{\BookRecordIndex{T0318}索伯列夫嵌入定理。} 有\BookRecordIndex{T0293}弱导数的函数不一定连续。例如在 $\Omega=B_1=\{x\in\mathbb R^n\mid |x|<1\}$ 上取 $u(x):=|x|^{-\alpha}$，其中 $\alpha\in\mathbb R$。此时 $\partial_j u=-\alpha x_j|x|^{-\alpha-2}$，归纳可得
\[
|\partial^\nu u(x)|\le c_\nu|x|^{-\alpha-|\nu|}.
\]
函数 $x\mapsto|x|^{-\beta}$ 在 $B_1$ 上可积，当且仅当 $\beta<n$。故只要 $\alpha p+kp<n$，$u$ 的不超过 $k$ 阶导数便都属于 $L^p$。若 $kp<n$，取 $0<\alpha<n/p-k$，便得到 $W^{k,p}(B_1)$ 中一个在原点不连续的函数。若 $kp>n$，这种构造失效；事实上此时每个 $W^{k,p}$ 函数都连续。

\begin{exercise}\label{ex:B-1-10}
设 $u:\mathbb R^n\setminus\{0\}\to\mathbb R$ 光滑，且在单位球外为零。假设 $u$ 及其一阶导数 $\partial_i u$ 都有有限的 $L^p$ 范数。证明 $u\in W^{1,p}(\mathbb R^n)$。提示：$u$ 及其一阶导数在原点附近不必有界；须证明\BookRecordIndex{T0293}弱导数与通常导数一致。为此，取检验函数 $\phi\in C_0^\infty(\mathbb R^n)$，考虑 $(\partial_i\phi)u$ 在 $\mathbb R^n\setminus B_\varepsilon$ 上的积分。利用下面的定理~\ref{thm:B-1-11} 或定理~\ref{thm:B-1-12}，证明 $\varepsilon\to0$ 时边界项趋于零。
\end{exercise}

对 $0<\mu\le1$，定义 Hölder 范数
\[
\|u\|_{C^{0,\mu}}=\sup_{\substack{x,y\in\Omega\\x\ne y}}\frac{|u(x)-u(y)|}{|x-y|^\mu}+\sup_{x\in\Omega}|u(x)|,
\qquad
\|u\|_{C^{k,\mu}}=\sum_{|\nu|\le k}\|\partial^\nu u\|_{C^{0,\mu}}.
\]
记 $C^{k,\mu}(\Omega)$ 为所有具有有限 Hölder 范数 $\|u\|_{C^{k,\mu}}$ 的 $C^k$ 函数 $u:\Omega\to\mathbb R$ 所成的空间。

\begin{theorem}\label{thm:B-1-11}
设 $\Omega\subset\mathbb R^n$ 是有界 \BookRecordIndex{T0035}Lipschitz 域，且 $kp>n$，$0<\mu:=k-n/p<1$。则存在常数 $c=c(k,p,\Omega)>0$，使对 $u\in C^\infty(\overline\Omega)$，
\[
\|u\|_{C^{0,\mu}}\le c\|u\|_{W^{k,p}}.
\]
包含映射 $W^{k,p}(\Omega)\hookrightarrow C^0(\overline\Omega)$ 是紧的。
\end{theorem}

\begin{theorem}\label{thm:B-1-12}
设 $\Omega\subset\mathbb R^n$ 是有界 \BookRecordIndex{T0035}Lipschitz 域，且 $kp<n$。则存在常数 $c=c(k,p,\Omega)>0$，使对 $u\in C^\infty(\overline\Omega)$，
\[
\|u\|_{L^{np/(n-kp)}}\le c\|u\|_{W^{k,p}}.
\]
若 $q<np/(n-kp)$，则包含映射 $W^{k,p}(\Omega)\hookrightarrow L^q(\Omega)$ 是紧的。
\end{theorem}

这些就是\emph{索伯列夫估计}。定理~\ref{thm:B-1-12} 的\BookRecordIndex{T0192}紧性结论称为
\emph{Rellich 定理}。由定理~\ref{thm:B-1-11}，若
$\Omega\subset\mathbb R^n$ 是有界 \BookRecordIndex{T0035}Lipschitz 域，则对 $1\le p\le\infty$，
\[
C^\infty(\overline\Omega)=\bigcap_{k=1}^\infty W^{k,p}(\Omega).
\]
$kp=n$ 是这些估计的临界情形。这时 $W^{k,p}$ 不嵌入连续函数空间。不过，对 $p'<p$ 应用定理~\ref{thm:B-1-12} 可知它嵌入每个 $1\le q<\infty$ 的 $L^q$。这里尤其关心 $n=p=2$、$k=1$ 的情形。

\begin{remark}\label{rem:B-1-13}
我们在引理~\ref{lem:10-4-1} 中见过，单位圆盘 $B_1=\{(x,y)\in\mathbb R^2\mid x^2+y^2<1\}$ 上存在一列函数 $u_j\in W^{1,2}(B_1)$，满足
\[
u_j(0)=1,\qquad\lim_{j\to\infty}\|u_j\|_{W^{1,2}}=0.
\]
二维情形下，$W^{1,2}$ 与 $C^0$ 之间的关系颇为微妙。例如，$f(e^{i\theta})=\theta$，$0\le\theta<2\pi$，这个有单一跳跃不连续点的函数，不能延拓为 $B_1$ 上的 $W^{1,2}$ 函数。另一方面，$S^1$ 上也有不连续函数可以延拓为 $B_1$ 上的 $W^{1,2}$ 函数。为证第一项，把 $f:S^1\to\mathbb R$ 展成 Fourier 级数
\[
f(e^{i\theta})=\sum_{k\in\mathbb Z}a_ke^{ik\theta},\qquad a_{-k}=\overline{a_k}.
\]
则 $f$ 能延拓为 $u\in W^{1,2}(B_1)$，当且仅当满足下列条件。\sourcefootnote{原书权重印为 $k$，但对实值函数有 $a_{-k}=\overline{a_k}$，双边和恒为零，不能给出所述迹空间判据；据此改为 $|k|$。}
\[
\sum_{k\in\mathbb Z}|k||a_k|^2<\infty.
\]
函数 $f(e^{i\theta})=\theta$ 不满足此条件。为证第二项，利用引理~\ref{lem:10-4-1} 证明中的论证构造无界函数 $u\in W^{1,2}(\mathbb R^2)$。
\end{remark}

\begin{exercise}\label{ex:B-1-14}
证明包含映射 $W^{1,2}(\mathbb R^n)\hookrightarrow L^2(\mathbb R^n)$ 不是紧的。
\end{exercise}

\begin{exercise}\label{ex:B-1-15}
本习题表明定理~\ref{thm:B-1-12} 中 \BookRecordIndex{T0035}Lipschitz 域的假设不能删去。考虑如下有界开集 $\Omega\subset\mathbb R^2$：
\[
\begin{split}
\Omega:={}&\{(x,y)\in\mathbb R^2\mid 0<x<1,\ 0<y<\tfrac12\}\\
&\cup\bigcup_{m=0}^\infty\{(x,y)\in\mathbb R^2\mid 2^{-2m-1}<x<2^{-2m},\ \tfrac12\le y<1\}.
\end{split}
\]
证明嵌入 $W^{1,2}(\Omega)\hookrightarrow L^2(\Omega)$ 不是紧的。找出一个光滑函数 $u\in W^{1,2}(\Omega)$，使它对任意 $q>2$ 都不属于 $L^q(\Omega)$。
\end{exercise}

定理~\ref{thm:B-1-11} 与定理~\ref{thm:B-1-12} 在 $k\ge2$ 时的断言容易由 $k=1$ 的情形\BookRecordIndex{T0326}推出。此外，鉴于命题~\ref{prop:B-1-9}，只须对 $W_0^{1,p}(\Omega)$ 证明这些结论。

\begin{lemma}\label{lem:B-1-16}
每个 $u\in C_0^\infty(\mathbb R^n)$ 都满足
\[
\sup_{x\ne y}\frac{|u(x)-u(y)|}{|x-y|^\mu}\le c\|\nabla u\|_{L^p},
\qquad\sup|u|\le c(\|u\|_{L^p}+\|\nabla u\|_{L^p}),
\]
其中 $p>n$，$\mu:=1-n/p$，$c:=2^{n+1}\omega_n^{-1/p}\bigl((p-1)/(p-n)\bigr)^{1-1/p}$，而 $\omega_n$ 表示单位球面 $S^{n-1}\subset\mathbb R^n$ 的面积。
\end{lemma}
\begin{proof}
先设 $B\subset\mathbb R^n$ 是内部非空的有界凸集，$u:B\to\mathbb R$ 是平均值为零的光滑函数。则
\begin{equation}\label{eq:B-1-3}
|u(x)|\le\frac{d^n\omega_n}{nV\omega_n^{1/p}}
\left(\frac{p-1}{p-n}\right)^{1-1/p}
d^{1-n/p}\|\nabla u\|_{L^p(B)},
\end{equation}
其中 $d:=\operatorname{diam}(B)$、$V:=\operatorname{Vol}(B)$。为证此式，先由 $\int_Bu=0$ 得
\[
u(x)=\frac1V\int_B\int_0^1
\langle\nabla u(x+t(y-x)),x-y\rangle\,dt\,dy.
\]
将 $\nabla u:B\to\mathbb R^n$ 在 $B$ 外定义为零，便有
\begin{align*}
V|u(x)|
&\le\int_{|y|\le d}\int_0^1|\nabla u(x+ty)|\,|y|\,dt\,dy\\
&=\int_0^d r^{n-1}\left(\int_{|\eta|=1}\int_0^1
|\nabla u(x+tr\eta)|\,r\,dt\,dS(\eta)\right)dr\\
&=\int_0^d r^{n-1}\left(\int_{|y|\le r}|y|^{1-n}|\nabla u(x+y)|\,dy\right)dr\\
&\le\frac{d^n}{n}\int_B|y-x|^{1-n}|\nabla u(y)|\,dy\\
&\le\frac{d^n}{n}\left(\int_{|y|\le d}|y|^{q-nq}\,dy\right)^{1/q}
\|\nabla u\|_{L^p(B)}.
\end{align*}
第二个等号中的 $dS(\eta)$ 表示球面的体积形式。第三个等号由换元 $dy=|y|^{n-1}r\,dt\,dS(\eta)$ 得出。最后一步是 $1/p+1/q=1$ 的 Hölder 不等式。计算积分即得~\eqref{eq:B-1-3}。现在取
\[
B=B_r(x_0),\qquad x_0:=\frac{x+y}{2},\qquad r:=\frac{|x-y|}{2}.
\]
这时 $d=|x-y|$，$d^n\omega_n/(nV)=2^n$。记 $u_B=V^{-1}\int_Bu$，于是
\begin{align*}
|u(x)-u(y)|
&\le|u(x)-u_B|+|u_B-u(y)|\\
&\le\frac{2^{n+1}}{\omega_n^{1/p}}
\left(\frac{p-1}{p-n}\right)^{1-1/p}
|x-y|^{1-n/p}\|\nabla u\|_{L^p(B)}
\end{align*}
对所有光滑 $u:\mathbb R^n\to\mathbb R$ 成立。这证明了引理的第一个断言。第二个不等式容易由~\eqref{eq:B-1-3} \BookRecordIndex{T0326}推出，只须取 $B=B_1(x)$ 并以 $u-u_B$ 代替 $u$。
\end{proof}

\begin{lemma}\label{lem:B-1-17}
设 $p<n$。则每个 $u\in C_0^\infty(\mathbb R^n)$ 满足
\[
\|u\|_{L^{np/(n-p)}}\le\frac{p}{\sqrt n}\frac{n-1}{n-p}\|\nabla u\|_{L^p}.
\]
\end{lemma}
\begin{proof}
下面这个优美的论证出自 Nirenberg。恒等式
\[
u(x)=\int_{-\infty}^{x_i}\partial_i u(x_1,\ldots,x_{i-1},t,x_{i+1},\ldots,x_n)\,dt
\]
表明
\[
|u(x)|^{n/(n-1)}\le\prod_{i=1}^n\left(\int_{-\infty}^\infty|\partial_i u(x)|\,dx_i\right)^{1/(n-1)}.
\]
注意右端虽然有 $n$ 个因子，但第 $i$ 个因子不依赖于 $x_i$。因此对 $x_i$ 积分时，只有 $n-1$ 个因子起作用。依次对 $x_1,\ldots,x_n$ 积分，每一步都用取 $m=n-1$ 的广义 Hölder 不等式
\[
\|v_1\cdots v_m\|_{L^1}\le\|v_1\|_{L^m}\cdots\|v_m\|_{L^m}.
\]
第 $k$ 步给出
\begin{align*}
\int|u|^{n/(n-1)}\,dx_1\cdots dx_k
\le{}&\prod_{i=1}^k\left(\int|\partial_i u|\,dx_1\cdots dx_k\right)^{1/(n-1)}\\
&\cdot\prod_{i=k+1}^n\left(\int|\partial_i u|\,dx_1\cdots dx_k\,dx_i\right)^{1/(n-1)}.
\end{align*}
这里第 $k+1$ 个因子不依赖 $x_{k+1}$。现对 $x_{k+1}$ 积分，便得到以 $k+1$ 代替 $k$ 的同一不等式。取 $k=n$，得到
\[
\|u\|_{L^{n/(n-1)}}\le
\prod_{i=1}^n\left(\int|\partial_i u|\right)^{1/n}
\le\frac1n\sum_{i=1}^n\int|\partial_i u|
\le\frac1{\sqrt n}\int|\nabla u|.
\]
最后一步利用向量 $a$ 的欧几里得范数 $|a|=\sqrt{\sum_i|a_i|^2}$。这证明了 $p=1$ 的情形。

为证一般情形，考虑函数 $v:=|u|^\alpha$ 的 $L^{n/(n-1)}$ 范数，其中 $\alpha:=p(n-1)/(n-p)$。由于 $|\nabla v|=\alpha|u|^{\alpha-1}|\nabla u|$ 且 $\alpha n/(n-1)=np/(n-p)$，得到
\begin{align*}
\left(\int|u|^{np/(n-p)}\right)^{1-1/n}
&\le\frac\alpha{\sqrt n}\int|u|^{\alpha-1}|\nabla u|\\
&\le\frac\alpha{\sqrt n}
\left(\int|u|^{\alpha q-q}\right)^{1/q}
\left(\int|\nabla u|^p\right)^{1/p}\\
&\le\frac\alpha{\sqrt n}
\left(\int|u|^{np/(n-p)}\right)^{1-1/p}
\left(\int|\nabla u|^p\right)^{1/p}.
\end{align*}
第二步是 $1/p+1/q=1$ 的 Hölder 不等式，最后一步使用 $\alpha q-q=np/(n-p)$。故引理对一般 $p$ 也成立。
\end{proof}

\begin{proof}[定理~\ref{thm:B-1-11} 与定理~\ref{thm:B-1-12} 的证明]
由引理~\ref{lem:B-1-16}，当 $p>n$ 时有嵌入 $W_0^{1,p}(\Omega)\hookrightarrow C^{0,1-n/p}(\overline\Omega)$；故由 Arzelà--Ascoli 定理，当 $\Omega$ 有界时，$W_0^{1,p}(\Omega)\hookrightarrow C^0(\overline\Omega)$ 是紧的。类似地，由引理~\ref{lem:B-1-17}，当 $p<n$、$q=np/(n-p)$ 时有嵌入 $W_0^{1,p}(\Omega)\hookrightarrow L^q(\Omega)$。要证明有界域 $\Omega$ 上 $q<np/(n-p)$ 时此嵌入的\BookRecordIndex{T0192}紧性，还须另作论证。不等式
\[
\|u\|_{L^q}\le\|u\|_{L^1}^\lambda
\|u\|_{L^{np/(n-p)}}^{1-\lambda},
\qquad\frac1q=\lambda+\frac{1-\lambda}{np/(n-p)}
\]
表明，只须证明有界域上的包含映射 $\iota:W_0^{1,p}(\Omega)\to L^1(\Omega)$ 是紧的。为此，记 $S_\delta:L^1(\Omega)\to L^1(\Omega)$ 为\BookRecordIndex{T0151}光滑化算子 $S_\delta f:=\rho_\delta*f$。由 Arzelà--Ascoli 定理，$S_\delta$ 是紧算子：若 $u_i$ 是 $L^1(\Omega)$ 中有界序列，则 $S_\delta u_i\in C^0(\overline\Omega)$ 一致有界且等度连续，因而有子序列在 $C^0(\overline\Omega)$ 中收敛，进而在 $L^1(\Omega)$ 中收敛。故复合映射 $S_\delta\circ\iota:W_0^{1,p}(\Omega)\to L^1(\Omega)$ 是紧的。

此外，积分下式
\begin{align*}
|u(x)-u_\delta(x)|
&=\left|\int_{|y|\le1}\rho(y)\int_0^\delta
\langle\nabla u(x-ty),y\rangle\,dt\,dy\right|\\
&\le\int_{|y|\le1}\rho(y)\int_0^\delta
|\nabla u(x-ty)|\,dt\,dy
\end{align*}
可得对 $u\in W_0^{1,p}(\Omega)$，
\[
\|u-S_\delta u\|_{L^1}\le\delta\|\nabla u\|_{L^1}
\le\delta\operatorname{Vol}(\Omega)^{1-1/p}\|\nabla u\|_{L^p}.
\]
这说明当 $\delta\to0$ 时，紧算子 $S_\delta\circ\iota$ 在一致算子拓扑下收敛到 $\iota$，故极限算子 $\iota$ 也紧。至此，已在以 $W_0^{1,p}(\Omega)$ 代替 $W^{k,p}(\Omega)$ 的情形下证明两个定理，且不要求 $\Omega$ 的边界具有任何\BookRecordIndex{T0408}正则性。要证明 $k=1$ 时的结论，只须把 $W_0^{1,p}$ 的相应嵌入定理与延拓定理（命题~\ref{prop:B-1-9}）结合。一般 $k$ 归约到 $k=1$ 的步骤留作练习。
\end{proof}

\paragraph{插值。} 为了直观理解\BookRecordIndex{T0319}索伯列夫空间，常把 $W^{k,q}$ 函数看作具有 $k-n/q$ 阶连续导数。这样，\BookRecordIndex{T0318}索伯列夫嵌入定理~\ref{thm:B-1-12} 可表述为：若 $W^{k,q}$ 函数比 $W^{j,p}$ 函数具有更多导数，即 $j\le k$ 且 $j-n/p\le k-n/q$，便有连续嵌入 $W^{k,q}\hookrightarrow W^{j,p}$。在临界情形 $k-n/q=0$ 时须谨慎。下述插值不等式的证明可见~\srcite{124}。

\begin{proposition}[Gagliardo--Nirenberg]\label{prop:B-1-18}
设 $\Omega\subset\mathbb R^n$ 是边界为 $C^k$ 类的有界开域。设 $j,k\ge0$ 为满足 $j<k$ 的整数，$1\le p,q,r\le\infty$，且 $k-n/q+n/r\ge0$，并有
\[
j-\frac np=\lambda\left(k-\frac nq\right)
+(1-\lambda)\left(-\frac nr\right),
\qquad\frac jk\le\lambda\le1.
\]
若 $(k-j)q=n$，再假设 $\lambda\ne1$。则存在常数 $c>0$，使对 $u\in W^{k,q}(\Omega)$，
\[
\|u\|_{W^{j,p}}\le c\|u\|_{W^{k,q}}^\lambda\|u\|_{L^r}^{1-\lambda}.
\]
\end{proposition}

\paragraph{乘积估计。} 应把 $kp>n$ 看成各种运算均能顺利进行的好情形，例如与光滑函数复合以及取乘积。

\begin{proposition}\label{prop:B-1-19}
设 $\Omega\subset\mathbb R^n$ 是边界为 $C^k$ 类的有界开域。假设 $kp>n$、$j\in\{1,\ldots,k\}$ 且 $j-n/q\le k-n/p$。则存在常数 $c>0$，使
\begin{align*}
\|uv\|_{W^{j,q}}&\le c\|u\|_{W^{j,q}}\|v\|_{W^{k,p}},\\
\|f\circ u\|_{W^{k,p}}&\le c\|f\|_{C^{k-1,1}}
\bigl(\|u\|_{L^\infty}^{k-1}+1\bigr)
\bigl(\|u\|_{W^{k,p}}+1\bigr)
\end{align*}
对 $u,v\in C^\infty(\Omega)$ 与 $f\in C^{k-1,1}(\mathbb R)$ 成立。
\end{proposition}

命题~\ref{prop:B-1-19} 及下一结果的证明是直接的练习；它们用到 Hölder 不等式及命题~\ref{prop:B-1-18} 中 Gagliardo--Nirenberg 的一般插值不等式。

\begin{proposition}\label{prop:B-1-20}
设 $\Omega\subset\mathbb R^n$ 是边界为 $C^k$ 类的有界开域。假设 $kp>n$，且 $f\in C^{\ell-1,1}(\mathbb R)$，其中 $\ell\ge k$。则映射
\[
W^{k,p}(\Omega)\to W^{k,p}(\Omega):u\mapsto f\circ u
\]
是巴拿赫空间之间的 $C^{\ell-k}$ 映射。
\end{proposition}

\paragraph{\BookRecordIndex{T0174}迹定理。} 理解具有\BookRecordIndex{T0293}弱导数的函数向低维子流形的限制并不容易。例如，$C^k$ 函数限制到超平面上仍为 $C^k$ 类这一显然事实，在\BookRecordIndex{T0319}索伯列夫空间中没有对应物。$W^{k,p}$ 函数限制到边界时会损失导数。

\begin{proposition}\label{prop:B-1-21}
设 $1<p<\infty$，$\Omega\subset\mathbb R^n$ 为边界光滑的有界开集。则存在 $c>0$，使每个 $u\in C^\infty(\overline\Omega)$ 满足
\[
\|u\|_{L^p(\partial\Omega)}
\le c\|u\|_{L^p(\Omega)}^{1-1/p}\|u\|_{W^{1,p}(\Omega)}^{1/p}.
\]
\end{proposition}
\begin{proof}
对 $x\in\partial\Omega$，记 $\nu(x)$ 为外法单位向量。取光滑函数 $f:\overline\Omega\to\mathbb R^n$，使 $f(x)=\nu(x)$ 对 $x\in\partial\Omega$ 成立。于是
\begin{align*}
\int_{\partial\Omega}|u|^p
&=\int_\Omega\operatorname{div}(f|u|^p)\\
&=\int_\Omega\bigl((\operatorname{div}f)|u|^p
+p\langle f,\nabla u\rangle u|u|^{p-2}\bigr)\\
&\le c_1\int_\Omega|u|^{p-1}(|u|+|\nabla u|)\\
&\le c_1\left(\int_\Omega|u|^p\right)^{(p-1)/p}
\left(\int_\Omega(|u|+|\nabla u|)^p\right)^{1/p}\\
&\le c_2\|u\|_{L^p}^{p-1}\|u\|_{W^{1,p}}.
\end{align*}
命题得证。
\end{proof}
\begin{proposition}\label{prop:B-1-22}
设 $p,\Omega$ 如命题~\ref{prop:B-1-21}，且 $u\in W^{k,p}(\Omega)$。则 $u\in W_0^{k,p}(\Omega)$ 当且仅当 $\partial^\alpha u$ 对 $|\alpha|\le k-1$ 在 $\partial\Omega$ 上为零。
\end{proposition}
\begin{proof}
若 $u\in W_0^{k,p}(\Omega)$，取 $C_0^\infty(\Omega)$ 中收敛到 $u$ 的序列 $u_i$，并对 $\partial^\alpha(u-u_i)$ 应用命题~\ref{prop:B-1-21}，便知 $|\alpha|\le k-1$ 时 $\partial^\alpha u$ 在 $\partial\Omega$ 上为零。证明逆命题时只须考虑 $k=1$。设 $u\in W^{1,p}(\Omega)$ 在 $\partial\Omega$ 上为零，并把 $u$ 在 $\mathbb R^n\setminus\Omega$ 上定义为零，从而延拓到整个 $\mathbb R^n$。该延拓属于 $W^{1,p}(\mathbb R^n)$。为见此，由命题~\ref{prop:B-1-4} 取 $\overline\Omega$ 上一列光滑函数 $u_j:\Omega\to\mathbb R$ 逼近 $u$。命题~\ref{prop:B-1-21} 给出 $u_j|_{\partial\Omega}\to0$ 于 $L^p(\partial\Omega)$。于是散度定理表明
\begin{align*}
\int_\Omega\bigl(u(\partial_i\phi)+(\partial_i u)\phi\bigr)
&=\lim_{j\to\infty}\int_\Omega
\bigl(u_j(\partial_i\phi)+(\partial_i u_j)\phi\bigr)\\
&=\lim_{j\to\infty}\int_{\partial\Omega}\nu^i u_j\phi=0
\end{align*}
对任意检验函数 $\phi\in C^\infty(\overline\Omega)$ 都成立，而不只是对 $\phi\in C_0^\infty(\Omega)$ 成立。因此延拓函数属于 $W^{1,p}(\mathbb R^n)$。由于 $u$ 在 $\Omega$ 外为零，现在可用一列在 $\partial\Omega$ 附近为零的 $W^{1,p}(\Omega)$ 函数逼近 $u$，这些函数因而属于 $W_0^{1,p}(\Omega)$。（练习！）命题得证。
\end{proof}

\paragraph{向量丛的截面。}

\begin{remark}\label{rem:B-1-23}
设 $M$ 是 $n$ 维光滑紧流形，$\pi:E\to M$ 是光滑向量丛。若截面 $s:M\to E$ 在所有局部坐标中的表示都属于 $W^{k,p}$，则称 $s$ 为 $W^{k,p}$ 类截面。此定义不依赖坐标的选择。为此，注意若 $\phi\in\operatorname{Diff}(\mathbb R^n)$，$\Phi\in C^\infty(\mathbb R^n,\mathbb R^{N\times N})$ 是矩阵值函数，且 $u\in W_{\mathrm{loc}}^{k,p}(\mathbb R^n,\mathbb R^N)$，则乘积 $\Phi\cdot(u\circ\phi)$ 对每个有界开集 $\Omega\subset\mathbb R^n$ 都属于 $W^{k,p}(\Omega)$，并有估计。\sourcefootnote{原书右端范数写作 $\|u\|_{W^{k,p}(\Omega)}$，但左端取 $u\circ\phi$，一般须在 $\phi(\Omega)$ 上控制 $u$；按复合映射定义域补正。原书 $\Phi(u\circ\phi)$ 依矩阵值 $\Phi$ 的类型在此理解为矩阵作用于向量，而非将 $u\circ\phi$ 代入 $\Phi$。}
\[
\|\Phi\cdot(u\circ\phi)\|_{W^{k,p}(\Omega)}
\le c\|u\|_{W^{k,p}(\phi(\Omega))},
\]
其中 $c=c(\phi,\Phi,\Omega)$ 与 $u$ 无关（见习题~\ref{ex:B-1-2}、\ref{ex:B-1-5}）。
即使 $kp\le n$，这一结论仍成立。要在 $W^{k,p}$ 截面空间上定义范数，可以取覆盖 $M$ 的有限个坐标卡上的 $W^{k,p}$ 范数之和。
\end{remark}

\begin{remark}\label{rem:B-1-24}
对取值于流形的函数 $u$，情形很不相同，因为这时必须考虑 $\phi\circ u$ 形式的复合，其中 $\phi$ 是坐标变换。命题~\ref{prop:B-1-19} 表明，对这样的函数需要 $kp>n$。更准确地，设 $X,M$ 是光滑闭流形，并设 $kp>n:=\dim X$。则空间
\[
\mathcal X^{k,p}:=W^{k,p}(X,M)
\]
可定义为由连续函数 $u:X\to M$ 组成的空间，它们在局部坐标卡中由 $W^{k,p}$ 函数表示。与注~\ref{rem:B-1-23} 一样，此定义不依赖坐标的选择。也可先选嵌入 $M\hookrightarrow\mathbb R^N$，再将 $\mathcal X^{k,p}$ 定义为 $W^{k,p}(X,\mathbb R^N)$ 中取值于 $M$ 的函数所成的子集。两种定义等价，留给读者验证。第二种定义也适用于 $kp\le n$：要求 $u\in W^{k,p}(X,\mathbb R^N)$ 几乎处处取值于 $M$。不过在这种情形下，空间 $\mathcal X^{k,p}$ 依赖于所选的嵌入。
\end{remark}
